% Generated draft: requires mathematical and source review.
\documentclass[11pt]{article}
\usepackage[margin=25mm]{geometry}
\usepackage{fontspec}
\setmainfont{Latin Modern Roman}
\usepackage{amsmath,amssymb,mathtools,unicode-math}
\setmathfont{Latin Modern Math}
\usepackage{enumitem,needspace,xurl,hyperref}
\setlength{\parindent}{0pt}
\setlength{\parskip}{5pt}
\setlength{\emergencystretch}{3em}
\newcommand{\N}{\mathbb{N}}
\newcommand{\Z}{\mathbb{Z}}
\newcommand{\Q}{\mathbb{Q}}
\newcommand{\R}{\mathbb{R}}
\newcommand{\C}{\mathbb{C}}
\newcommand{\F}{\mathbb{F}}
\newcommand{\A}{\mathbb{A}}
\newcommand{\PP}{\mathbb{P}}
\newcommand{\E}{\mathbb{E}}
\newcommand{\eps}{\varepsilon}
\newcommand{\dd}{\,\mathrm{d}}
\begin{document}
\section*{UnsolvedMath ID: 30003845}
\textbf{Source problem code:} catalogue entry
\subsection{Definitions and mathematical statement}
Put \(I=[0,1]\). Use the standard convention that \(\operatorname{PL}_o(I)\) is the group, under composition, of increasing piecewise-linear homeomorphisms \(I\to I\) with finitely many linear pieces. In particular, these homeomorphisms fix both endpoints. Thompson's group \(F\) is the subgroup consisting of maps whose breakpoints are dyadic rationals and whose slopes belong to \(\{2^k:k\in\Z\}\).

For groups \(G,H\), write \(G\preccurlyeq H\) if there exists an injective group homomorphism \(G\to H\). Say that \(G\) and \(H\) are biembeddable, and write \(G\sim H\), if both \(G\preccurlyeq H\) and \(H\preccurlyeq G\). For an ambient group \(A\), let
\[
\mathcal B(A)=\{[H]_{\sim}:H\leq A\},
\qquad
[G]_{\sim}\leq[H]_{\sim}\ \Longleftrightarrow\ G\preccurlyeq H.
\]
This defines a partial order on the biembeddability classes. A partially ordered set \((P,\leq)\) is well-quasi-ordered if every sequence \((p_n)_{n\in\N}\) in \(P\) has indices \(i<j\) with \(p_i\leq p_j\).

\textbf{Questions.}
\begin{enumerate}
\item Is \(\mathcal B(F)\) well-quasi-ordered? Equivalently, is it true that
\[
\forall (H_n)_{n\in\N}\in\{H:H\leq F\}^{\N},\quad
\exists i<j\quad H_i\preccurlyeq H_j?
\]
\item Does \(\mathcal B(\operatorname{PL}_o(I))\) fail to be well-quasi-ordered? Equivalently, does there exist a sequence \((K_n)_{n\in\N}\) of subgroups of \(\operatorname{PL}_o(I)\) such that
\[
\forall i<j,\quad K_i\not\preccurlyeq K_j?
\]
\end{enumerate}

\textbf{Scope and unresolved qualifications.} The supplied statement says subgroups without imposing finite generation, countability, or other subgroup restrictions; the displayed questions retain that unrestricted scope. Embeddings are interpreted here as abstract injective group homomorphisms, without an additional condition relating the actions on \(I\). The source record does not explicitly define its embedding convention or \(\operatorname{PL}_o(I)\); the standard conventions above require confirmation against the cited report. The second question is interpreted as asking for a family of subgroup classes witnessing failure of well-quasi-ordering. A single class alone cannot witness such failure.
\subsection{Short English statement}
Does every sequence of subgroups of Thompson's group \(F\) contain an earlier group that embeds into a later one? Can a sequence of subgroups of \(\operatorname{PL}_o([0,1])\) fail this property? The draft uses abstract group embeddings and retains the source's unrestricted subgroup quantifier. The supplied open status is dated curation, not an independently verified conclusion.
\subsection{Sources}
\begin{itemize}
\item Supplied local UnsolvedMath registry record, permanent ID \(30003845\), title \emph{Well-Quasi-Ordering of Thompson Subgroups}, exact source problem code \texttt{catalogue entry}, in \emph{Oberwolfach Reports Open Problems}. The record lists proposed year 2018 and supplies no individual proposer.
\item Supplied source citation: \emph{Cohomological and Metric Properties of Groups of Homeomorphisms of R} (2019), Oberwolfach Reports, pp.~1579--1633, DOI \texttt{10.4171/owr/2018/26}. \url{https://doi.org/10.4171/owr/2018/26}. The background attributes the collection to workshop open-problem discussions documented in Oberwolfach Reports.
\item Supplied literature reference: Bleak--Brin--Moore, \url{https://doi.org/10.4171/JCA/49}. The supplied assessment attributes to these authors a long well-ordered restricted family and incomparable examples; it does not specify that family's hypotheses or order type.
\item Supplied literature reference labeled \emph{Oberwolfach Report 2019}: \url{https://ems.press/content/serial-article-files/46748}.
\item Literature assessment date: \texttt{2026-08-21}. This is supplied dated curation, not independently verified. The supplied assessment labels both the full biembeddability well-quasi-ordering question and the \(\operatorname{PL}_o(I)\) violation question open. The background expressly requests independent verification of time-sensitive status and cited claims.
\item The registry labels the statement \texttt{corrected\_verified} and reports a manual semantic, typographic, and LaTeX review. These are supplied editorial assertions, not independent mathematical or source verification for this draft.
\end{itemize}
\end{document}
